% Generated by scripts/blueprint_gen.py; do not edit.
\chapter{The triangulated structure on Ho(P\textsubscript{A})}\label{chap:Z2}
\textbf{Targets.}

\begin{itemize}
\item \textbf{Z2.} For every bounded graded-commutative curved dga \texttt{𝒜} (the project's \texttt{GCDga}, curvature \texttt{c} allowed), \texttt{HoP 𝒜} gets:
\begin{itemize}
\item a shift \texttt{HasShift (HoP 𝒜) ℤ};
\item a zero object;
\item \texttt{Pretriangulated (HoP 𝒜)}, with distinguished triangles those isomorphic to cone triangles;
\item \texttt{IsTriangulated (HoP 𝒜)} (octahedron).
This is [B, Prop 2.4.1]: "P\textsubscript{A} is pretriangulated in the sense of Bondal–Kapranov, so Ho P\textsubscript{A} is triangulated". Block's proof is the single line "same as Prop 1 and Prop 2 of [BK]".

\end{itemize}
\item \textbf{K4} (blueprint §9, original form). For the unit module \texttt{𝒜} over the Dolbeault algebra (\texttt{Q = 𝒜⁰} in degree 0, \texttt{𝕏 = ∂̄}): \texttt{Hom\_\{HoP\}(𝒜, 𝒜⟦k⟧) ≅ H\textasciicircum{}\{0,k\}\_∂̄(X)} as ℂ-vector spaces, for every \texttt{k : ℤ}. Both sides vanish for \texttt{k < 0} and \texttt{k > dim X}.
\end{itemize}
\textbf{Route} (research of 1 Nov 2026). Direct construction on \texttt{HoP}, imitating Mathlib's \texttt{HomotopyCategory} files: \texttt{Shift} (320 lines), \texttt{MappingCone} (671), \texttt{Pretriangulated} (534) and \texttt{Triangulated} (229). The other options fail:

\begin{itemize}
\item Mathlib has no dg-categories or dg-modules.
\item The total-complex functor is not full, and curvature rules out ordinary complexes.
\item Transport along \texttt{J : HoP ≅ D\_perf} only works for the Dolbeault algebra, needs the main theorem, and gives a non-canonical shift. K4 would then not be a check of the shift.
\end{itemize}
Membership in \texttt{Hom\textasciicircum{}n}: \texttt{φ(α·x) = (−1)\textasciicircum{}\{n|α|\} α·φ(x)}. Differential: \texttt{δφ = 𝕏 ∘ φ − (−1)\textasciicircum{}n φ ∘ 𝕏} (A4, equal to Mathlib's \texttt{HomComplex.δ} by A4.d).

\begin{itemize}
\item \textbf{Shift.}
\begin{itemize}
\item \texttt{Q[a]\textasciicircum{}p := Q\textasciicircum{}\{p+a\}}.
\item \texttt{ι\_a : Tot(Q[a]) → Tot(Q)}, \texttt{α ⊗ e ↦ (−1)\textasciicircum{}\{a|α|\} α ⊗ e} (\texttt{GCDga.twist a} ⊗ reindexing by \texttt{Equiv.addRight a}). It is bijective and lies in \texttt{Hom\textasciicircum{}\{a\}}: an element of degree \texttt{m} goes to degree \texttt{m + a}, and \texttt{ι\_a(β x) = (−1)\textasciicircum{}\{a|β|\} β ι\_a(x)}.
\item \texttt{𝕏[a] := (−1)\textasciicircum{}a ι\_a⁻¹ ∘ 𝕏 ∘ ι\_a}. Checks:
\begin{itemize}
\item Leibniz: \texttt{ι\_a⁻¹(dβ · y) = (−1)\textasciicircum{}\{a(|β|+1)\} dβ ι\_a⁻¹ y}; the signs cancel to \texttt{𝕏[a](βx) = dβ·x + (−1)\textasciicircum{}\{|β|\} β 𝕏[a]x}.
\item Flatness: \texttt{𝕏[a]² = ι⁻¹ c ι = c}, since \texttt{c} has even degree.
\item On generators: \texttt{𝕏[a]\textasciicircum{}j = (−1)\textasciicircum{}\{a(j+1)\} 𝕏\textasciicircum{}j}, consistent with Mathlib's \texttt{n.negOnePow} on differentials.
\end{itemize}
\item Generic lemma (Z2.a): if \texttt{u : Tot' → Tot} is bijective, \texttt{u ∈ Hom\textasciicircum{}m} (semilinear of degree \texttt{m}), and \texttt{𝕏} is a cohesive operator, then \texttt{(−1)\textasciicircum{}m u⁻¹ 𝕏 u} is one. This is cleaner than going through \texttt{GenData}.
\item Morphisms: \texttt{φ ↦ ι\_a⁻¹ φ ι\_a}, with \texttt{δ(ι⁻¹hι) = (−1)\textasciicircum{}a ι⁻¹(δh)ι}, so boundaries go to boundaries.
\item \texttt{M[a][b] ≅ M[a+b]} is the identity on generators up to reindexing; the signs agree (\texttt{(−1)\textasciicircum{}\{(a+b)(j+1)\}}).
\end{itemize}
\item **Cone of \texttt{φ ∈ Z0(M, N)}.**
\begin{itemize}
\item \texttt{Q\_C\textasciicircum{}p := Q\_N\textasciicircum{}p × Q\_M\textasciicircum{}\{p+1\}}.
\item Structure maps, all A-semilinear and defined on generators via \texttt{homEquivGen} (A4.c): \texttt{inl ∈ Hom\textasciicircum{}\{−1\}(M, C)}, \texttt{inr ∈ Hom\textasciicircum{}0(N, C)}, \texttt{fst ∈ Hom\textasciicircum{}1(C, M)}, \texttt{snd ∈ Hom\textasciicircum{}0(C, N)}.
\item Identities: \texttt{fst∘inl = id}, \texttt{snd∘inr = id}, \texttt{fst∘inr = 0}, \texttt{snd∘inl = 0}, \texttt{inl∘fst + inr∘snd = id}.
\item Operator: \texttt{𝕏\_C := inr 𝕏\_N snd + inr φ fst − inl 𝕏\_M fst}. On \texttt{inr f + inl e} this is Block's matrix. Flatness from those of \texttt{M}, \texttt{N} and \texttt{δφ = 0}.
\item Then \texttt{δ(inl) = inr ∘ φ}, \texttt{δ(inr) = 0}, \texttt{δ(fst) = 0}, \texttt{δ(snd) = −φ ∘ fst} (Mathlib's \texttt{δ\_inl}, \texttt{δ\_snd}).
\end{itemize}
\item \textbf{Triangle.} \texttt{M →φ N →inr C(φ) →(−ι₁⁻¹∘fst) M⟦1⟧}, where \texttt{ι₁⁻¹ ∘ fst ∈ Z0(C, M[1])}. This is Mathlib's \texttt{Cocycle.homOf ((-fst φ).rightShift 1 0 \_)}.
\end{itemize}

\begin{theorem}[Z2.a — Shift]
\label{Z2.a}
\lean{Cohesive.shiftQ, Cohesive.shiftIota, Cohesive.QuasiCohesive.transport, Cohesive.CohesiveModule.shift, Cohesive.shiftIota_mem_hom, Cohesive.δ_shiftIota}
\leanok
Let A be a bounded graded-commutative curved dga over k and M = (Q, 𝕏) a cohesive module, a ∈ ℤ. The shifted generators are Q[a]\textsuperscript{p} = Q\textsuperscript{p+a}, and ι\textsubscript{a} : Tot Q[a] → Tot Q, α ⊗ e ↦ (−1)\textsuperscript{a|α|} α ⊗ e, is a k-linear bijection taking total degree m to m + a with ι\textsubscript{a}(βx) = (−1)\textsuperscript{a|β|} β ι\textsubscript{a}(x).

\begin{itemize}
\item Generally, if u : Tot Q′ → Tot Q is such a bijection of degree m, then (−1)\textsuperscript{m} u\textsuperscript{−1} 𝕏 u is a quasi-cohesive structure on Q′.
\item The shift M[a] = (Q[a], (−1)\textsuperscript{a} ι\textsubscript{a}\textsuperscript{−1} 𝕏 ι\textsubscript{a}) is a cohesive module, and ι\textsubscript{a} ∈ Hom\textsuperscript{a}(M[a], M) is closed.
\end{itemize}
Deviation: Block's 𝔼[1] = −𝔼 fails the Leibniz rule for left modules; the shift is transported along ι\textsubscript{a} instead.
\par\smallskip{\footnotesize Source: blueprint-tr.md §5.}
\end{theorem}
\begin{proof}
\leanok
\uses{A4}
See the Lean proof and the cited items.
\paragraph{Blueprint notes.}
\texttt{ι\_a}, generic transport lemma, \texttt{CohesiveModule.shift}

Files: \texttt{Algebra/Shift.lean}.
\end{proof}

\begin{theorem}[Z2.b — Shift functor on Ho(P\textsubscript{A})]
\label{Z2.b}
\lean{Cohesive.shiftMkCoreHoP, Cohesive.isShiftData_shiftIota, Cohesive.IsShiftData.trans, Cohesive.shiftCompZ0, Cohesive.shiftCompIso, Cohesive.conjL}
\leanok
Let A be a bounded graded-commutative curved dga over k. A \emph{shift datum} of degree m from M′ to M is a k-linear bijection v : Tot M′ → Tot M taking total degree n onto n + m, with v(βx) = (−1)\textsuperscript{m|β|} β v(x) and 𝕏′ = (−1)\textsuperscript{m} v\textsuperscript{−1} 𝕏 v; shift data compose (degrees add), and two shift data v₁, v₂ of the same degree onto M give an isomorphism v₂\textsuperscript{−1} ∘ v₁ in Ho(P\textsubscript{A}).

\begin{itemize}
\item Ho(P\textsubscript{A}) has a shift by ℤ: M⟦a⟧ = M[a] on objects, [φ]⟦a⟧ = [ι\textsubscript{a}\textsuperscript{−1} ∘ φ ∘ ι\textsubscript{a}] on morphisms, with M[0] ≅ M and M[a+b] ≅ M[a][b] the comparisons of the shift data ι\textsubscript{0}, id and ι\textsubscript{a+b}, ι\textsubscript{a} ∘ ι\textsubscript{b}.
\item Each shift functor is additive and k-linear.
\end{itemize}
\par\smallskip{\footnotesize Source: blueprint-tr.md §5.}
\end{theorem}
\begin{proof}
\leanok
\uses{A5, Z2.a}
See the Lean proof and the cited items.
\paragraph{Blueprint notes.}
shift on \texttt{Hom}/\texttt{Z0}/\texttt{H0}, \texttt{HasShift (HoP 𝒜) ℤ} (\texttt{hasShiftMk}), additivity, \texttt{k}-linearity

Files: \texttt{Algebra/ShiftHoP.lean}.
\end{proof}

\begin{theorem}[Z2.c — Zero object]
\label{Z2.c}
\lean{Cohesive.prodQ, Cohesive.dsumProdEquiv, Cohesive.zeroCohesive, Cohesive.zeroHoP, Cohesive.isZero_zeroHoP}
\leanok
Let A be a bounded graded-commutative curved dga over k.

\begin{itemize}
\item For families Q₁, Q₂ of A\textsuperscript{0}-modules, with (Q₁ × Q₂)\textsuperscript{p} = Q₁\textsuperscript{p} × Q₂\textsuperscript{p}, the map α ⊗ (e₁, e₂) ↦ (α ⊗ e₁, α ⊗ e₂) is an A-linear isomorphism Tot(Q₁ × Q₂) ≅ Tot Q₁ × Tot Q₂ preserving total degree componentwise.
\item The module with all Q\textsuperscript{p} = 0 and 𝕏 = 0 is cohesive and is a zero object of Ho(P\textsubscript{A}); in particular Ho(P\textsubscript{A}) has a zero object.
\end{itemize}
\par\smallskip{\footnotesize Source: blueprint-tr.md §5.}
\end{theorem}
\begin{proof}
\leanok
\uses{A5}
See the Lean proof and the cited items.
\paragraph{Blueprint notes.}
\texttt{Tot.prodEquiv} (v1.18, below), zero object

Files: \texttt{Algebra/TotProd.lean}, \texttt{Algebra/HoPZero.lean}.
\end{proof}

\begin{definition}[Z2.d — Cone]
\label{Z2.d}
\lean{Cohesive.coneQ, Cohesive.cone, Cohesive.coneX}
\leanok
\uses{Z2.a, Z2.c}
Let A be a bounded graded-commutative curved dga over k, M, N cohesive modules and φ ∈ Z\textsuperscript{0}(M, N) a closed degree-0 map. The cone C(φ) has generators Q\textsubscript{N}\textsuperscript{p} × Q\textsubscript{M}\textsuperscript{p+1}; through Z2.c and ι\textsubscript{1}, Tot C(φ) ≅ Tot M × Tot N, giving inl ∈ Hom\textsuperscript{−1}(M, C), inr ∈ Hom\textsuperscript{0}(N, C), fst ∈ Hom\textsuperscript{1}(C, M), snd ∈ Hom\textsuperscript{0}(C, N) with fst∘inl = id, snd∘inr = id, fst∘inr = 0, snd∘inl = 0, inl∘fst + inr∘snd = id. The operator is 𝕏\textsubscript{C} = inr 𝕏\textsubscript{N} snd + inr φ fst − inl 𝕏\textsubscript{M} fst, and C(φ) is a cohesive module with δ inl = inr ∘ φ, δ inr = 0, δ fst = 0, δ snd = −φ ∘ fst (Mathlib's mapping-cone signs).
\par\smallskip{\footnotesize Source: blueprint-tr.md §5.}
\end{definition}
\paragraph{Blueprint notes.}
cone, structure maps, identities, \texttt{δ} formulas

Files: \texttt{Algebra/Cone.lean}.

\begin{definition}[Z2.e — Maps of cones]
\label{Z2.e}
\lean{Cohesive.coneInrZ0, Cohesive.coneDesc, Cohesive.coneLift, Cohesive.coneMap}
\leanok
\uses{Z2.d}
In the setting of Z2.d (φ ∈ Z\textsuperscript{0}(M, N)), with T a cohesive module:

\begin{itemize}
\item inr : N → C(φ) is closed of degree 0;
\item (desc) for α ∈ Hom\textsuperscript{−1}(M, T) and closed β : N → T of degree 0 with δα = β ∘ φ, α ∘ fst + β ∘ snd : C(φ) → T is closed of degree 0;
\item (lift) for closed α ∈ Hom\textsuperscript{1}(T, M) and β ∈ Hom\textsuperscript{0}(T, N) with δβ = −φ ∘ α, inl ∘ α + inr ∘ β : T → C(φ) is closed of degree 0;
\item (map of cones) for φ\textsubscript{i} ∈ Z\textsuperscript{0}(M\textsubscript{i}, N\textsubscript{i}), closed a : M₁ → M₂, b : N₁ → N₂ of degree 0 and h ∈ Hom\textsuperscript{−1}(M₁, N₂) with δh = b ∘ φ₁ − φ₂ ∘ a, desc(inl₂ ∘ a + inr₂ ∘ h, inr₂ ∘ b) is a closed degree-0 map C(φ₁) → C(φ₂).
\end{itemize}
\par\smallskip{\footnotesize Source: blueprint-tr.md §5.}
\end{definition}
\paragraph{Blueprint notes.}
\texttt{lift}/\texttt{desc}, maps of cones from homotopies

Files: \texttt{Algebra/ConeMaps.lean}.

\begin{theorem}[Z2.f — Ho(P\textsubscript{A}) pretriangulated]
\label{Z2.f}
\lean{Cohesive.distinguishedTriangles, Cohesive.contractible_distinguished, Cohesive.distinguished_cocone_triangle, Cohesive.exists_coneTriangle_hom, Cohesive.complete_distinguished_triangle_morphism, Cohesive.rotateTriangleIso}
\leanok
Let A be a bounded graded-commutative curved dga over k (no further hypotheses). For φ ∈ Z\textsuperscript{0}(M, N) the \emph{cone triangle} is M →[φ] N →[inr] C(φ) →[w\textsubscript{φ}] M⟦1⟧ with w\textsubscript{φ} = −ι\textsubscript{1}\textsuperscript{−1} ∘ fst, and a triangle in Ho(P\textsubscript{A}) is \emph{distinguished} if it is isomorphic to a cone triangle. With the shift of Z2.b, this makes Ho(P\textsubscript{A}) a pretriangulated category:

\begin{itemize}
\item X →[id] X → 0 → X⟦1⟧ is distinguished, and every morphism is the first map of a distinguished triangle;
\item a commutative square between the first maps of two distinguished triangles extends to a morphism of triangles (for cone triangles, by a map of cones, Z2.e);
\item (cone triangle of φ).rotate ≅ cone triangle of inr\textsubscript{φ}, with components id, id and a homotopy equivalence M⟦1⟧ ≅ C(inr\textsubscript{φ}); distinguished triangles are closed under rotation.
\end{itemize}
\par\smallskip{\footnotesize Source: blueprint-tr.md §5.}
\end{theorem}
\begin{proof}
\leanok
\uses{Z2.a, Z2.b, Z2.e}
Notation: \texttt{ι\_a : Tot(M[a]) → Tot(M)} (Z2.a, \texttt{Hom\textasciicircum{}a}, closed), \texttt{w\_φ := −ι₁⁻¹ ∘ fst ∈ Z⁰(C(φ), M[1])} (P14: Mathlib's \texttt{triangle φ}.mor₃ with no extra sign on generators), \texttt{f[a] := ι\_a⁻¹ f ι\_a} (\texttt{conjZ0}, the shift on morphisms, Z2.b). Primes denote the structure maps of \texttt{C(inr\_φ)}. Mathlib's \texttt{leftShift} signs (\texttt{(a n' + a(a−1)/2)}) are absorbed into these formulas.

\begin{itemize}
\item \textbf{Triangle and set} (\texttt{Algebra/ConeTriangle.lean}, Z2.f1). \texttt{coneTriangle φ := (M →[φ] N →[inr] C(φ) →[w\_φ] M⟦1⟧)}; distinguished = isomorphic to some \texttt{coneTriangle φ} (\texttt{φ ∈ Z⁰}, representatives in \texttt{HoP}).
\begin{itemize}
\item contractible: \texttt{C(id\_M)} is zero: \texttt{δ(inl ∘ snd) = id} (\texttt{δ inl ∘ snd − inl ∘ δ snd = inr snd + inl fst}), \texttt{isZero\_of\_contraction};
\item cocone: choose a representative of \texttt{f};
\item morphisms: a square commuting in \texttt{HoP} gives \texttt{h ∈ Hom\textasciicircum{}\{−1\}} with \texttt{δh = bφ₁ − φ₂a} (\texttt{B⁰}); \texttt{coneMap} (Z2.e); \texttt{a[1] ∘ w₁ = w₂ ∘ coneMap} exactly (\texttt{fst₂ ∘ coneMap = a ∘ fst₁}).
\end{itemize}
\item \textbf{Rotation} (\texttt{Algebra/ConeRotate.lean}, Z2.f2; Mathlib \texttt{rotateHomotopyEquiv}). \texttt{hom := coneLift inr\_φ (φ ∘ ι₁) (−inl ∘ ι₁) : M[1] → C(inr\_φ)} (\texttt{δ(−inl ι₁) = −inr φ ι₁}), \texttt{inv := coneDesc inr\_φ 0 w\_φ : C(inr\_φ) → M[1]}; \texttt{inv ∘ hom = id} exactly, \texttt{hom ∘ inv − id = δ H} with \texttt{H := −inl' ∘ snd\_φ ∘ snd'}. Triangle iso \texttt{(coneTriangle φ).rotate ≅ coneTriangle inr\_φ} with components \texttt{(id, id, [hom])}: \texttt{w' ∘ hom = −φ[1]} exactly (\texttt{fst' ∘ hom = φ ι₁}), and \texttt{hom ∘ w\_φ − inr' = δ H₂}, \texttt{H₂ := −inl' ∘ snd\_φ}.
\item \textbf{Shift} (\texttt{Algebra/ConeShift.lean}, Z2.f3; Mathlib \texttt{shiftIso}, \texttt{shiftTriangleIso}). \texttt{C(φ)[n] ≅ C(φ[n])}: \texttt{hom := coneLift φ[n] ((−1)\textasciicircum{}n ι\_n⁻¹ fst ι\_n) (ι\_n⁻¹ snd ι\_n)}, \texttt{inv := coneDesc φ[n] ((−1)\textasciicircum{}n ι\_n⁻¹ inl ι\_n) (ι\_n⁻¹ inr ι\_n)}, inverse exactly; triangle iso \texttt{(Triangle.shiftFunctor n).obj (coneTriangle φ) ≅ coneTriangle φ[n]} with components \texttt{(id, (−1)\textasciicircum{}n, shiftIso)}; the third square involves \texttt{shiftFunctorComm} of \texttt{HoP} (Z2.b's \texttt{shiftFunctorAdd}), the main risk.
\item \textbf{Instance} (\texttt{Algebra/Pretriangulated.lean}): as Mathlib, \texttt{invRotate} via \texttt{invRotateIsoRotateRotateShiftFunctorNegOne} and shift-closure.
\end{itemize}
\paragraph{Blueprint notes.}
the pretriangulated axioms, shift of cones

Files: \texttt{Algebra/Pretriangulated.lean}.
\end{proof}

\begin{theorem}[Z2.g — Octahedron axiom]
\label{Z2.g}
\lean{Cohesive.compTriangle, Cohesive.compTriangle_distinguished, Cohesive.instIsTriangulatedHoP}
\leanok
Let A be a bounded graded-commutative curved dga over k. For closed degree-0 maps u : X₁ → X₂, v : X₂ → X₃, the triangle C(u) →[m₁] C(v ∘ u) →[m₂] C(v) →[w\textsubscript{v} ≫ inr\textsubscript{u}⟦1⟧] C(u)⟦1⟧, with m₁, m₂ the maps of cones of the squares (id, v) and (u, id) with zero homotopy, is distinguished. Hence the pretriangulated category Ho(P\textsubscript{A}) satisfies the octahedron axiom: Ho(P\textsubscript{A}) is triangulated, for every curvature c.
\par\smallskip{\footnotesize Source: blueprint-tr.md §5.}
\end{theorem}
\begin{proof}
\leanok
\uses{Z2.a, Z2.b, Z2.c, Z2.d, Z2.e, Z2.f}
Mathlib \texttt{HomotopyCategory/Triangulated.lean} (\texttt{mappingConeCompTriangle}, \texttt{mappingConeCompHomotopyEquiv}, \texttt{IsTriangulated.mk'}). For \texttt{u ∈ Z⁰(X₁, X₂)}, \texttt{v ∈ Z⁰(X₂, X₃)}, \texttt{uv := v ∘ u}. Conventions as in §3c (\texttt{δ inl = inr φ}, \texttt{δ snd = −φ fst}, \texttt{δ fst = δ inr = 0}, \texttt{δ(f∘g) = δf∘g + (−1)\textasciicircum{}\{|f|\} f∘δg}).

\begin{itemize}
\item \textbf{Comparison maps} (\texttt{Algebra/ConeComp.lean}). \texttt{m₁ := coneMap u uv id v 0 : C(u) → C(uv)}, \texttt{m₂ := coneMap uv v u id 0 : C(uv) → C(v)} (homotopy \texttt{0}: \texttt{v∘u − uv∘id = 0}, \texttt{id∘uv − v∘u = 0}). Triangle \texttt{C(u) →[m₁] C(uv) →[m₂] C(v) →[w\_v ≫ inr\_u⟦1⟧] C(u)⟦1⟧}.
\item \textbf{Homotopy equivalence} \texttt{C(v) ≃ C(m₁) =: D}. \texttt{hom := coneLift m₁ (inr\_u ∘ fst\_v) (inr\_\{uv\} ∘ snd\_v)} (\texttt{δ(inr\_\{uv\} snd\_v) = −inr\_\{uv\} v fst\_v = −m₁ inr\_u fst\_v} by \texttt{m₁ inr\_u = inr\_\{uv\} v}); \texttt{inv := coneDesc m₁ (inl\_v ∘ snd\_u) m₂} (\texttt{δ(inl\_v snd\_u) = inr\_v v snd\_u + inl\_v u fst\_u = m₂ ∘ m₁}, using \texttt{m₁ = inl\_\{uv\} fst\_u + inr\_\{uv\} v snd\_u}). \texttt{inv ∘ hom = id} exactly (\texttt{fst\_v inl\_v = id}, \texttt{snd\_v m₂ = snd\_\{uv\}}); \texttt{hom ∘ inv − id = δH} with \texttt{H := −inl\_D ∘ inl\_u ∘ fst\_\{uv\} ∘ snd\_D} (checked by hand: with \texttt{K = inl\_u fst\_\{uv\} snd\_D}, \texttt{δ(inl\_D K) = inr\_D inl\_\{uv\} fst\_\{uv\} snd\_D − inl\_D inr\_u u fst\_\{uv\} snd\_D + inl\_D inl\_u fst\_u fst\_D}, which is \texttt{−(hom ∘ inv − id)} after expanding \texttt{id\_D}, \texttt{id\_\{C(u)\}}, \texttt{id\_\{C(uv)\}}).
\item \textbf{Distinguished.} Triangle iso \texttt{(id, id, [hom])} to \texttt{coneTriangle m₁}: \texttt{inv ∘ inr\_D = m₂} exactly (so \texttt{m₂ ≫ [hom] = [inr\_D]} in \texttt{HoP}), and \texttt{w\_D ∘ hom = −ι₁⁻¹ inr\_u fst\_v = inr\_u[1] ∘ w\_v} exactly.
\item \textbf{Octahedron} (\texttt{Algebra/Triangulated.lean}). \texttt{IsTriangulated.mk'} with the three cone triangles of \texttt{u}, \texttt{v}, \texttt{uv}, \texttt{Iso.refl}s, and \texttt{Octahedron.mk [m₁] [m₂]}: \texttt{m₁ inr\_u = inr\_\{uv\} v}, \texttt{w\_\{uv\} m₁ = w\_u} (\texttt{fst m₁ = fst}), \texttt{m₂ inr\_\{uv\} = inr\_v}, \texttt{u[1] w\_\{uv\} = w\_v m₂} (\texttt{fst\_v m₂ = u fst\_\{uv\}}), all exact; the last field is the distinguished triangle above.
\end{itemize}
\paragraph{4. K4}
\begin{itemize}
\item \textbf{(i) Generic} (model: Mathlib \texttt{HomComplexCohomology.homAddEquiv}). \texttt{Hom\_\{HoP\}(M, N⟦k⟧) ≅ H\textasciicircum{}k(Hom\textasciicircum{}•(M,N))} via \texttt{φ ↦ ι\_k ∘ φ}, since \texttt{δ(ι\_k ∘ φ) = ± ι\_k ∘ δφ}. The cohomology is the quotient \texttt{(Hom\textasciicircum{}k ⊓ ker δ) ⧸ δ(Hom\textasciicircum{}\{k−1\})} (like \texttt{H0}), or the homology of A6's \texttt{homComplex}.
\item \textbf{(ii) Unit module.} \texttt{Hom\textasciicircum{}k(U, U) ≅ 𝒜\textasciicircum{}k} by evaluation at \texttt{1 ⊗ gen} (\texttt{homEquivGen}, \texttt{hom\_tmul}), with \texttt{δ ↔ d} and no sign (\texttt{θ(gen) = 0}). \texttt{unitCohesive} exists in \texttt{ChecksAlgebra.lean:338}; move it to a library file (it is not frozen).
\item \textbf{(iii) Dolbeault.} Define \texttt{H\textasciicircum{}k(𝒜, d) := (𝒜\textasciicircum{}k ⊓ ker d) ⧸ d(𝒜\textasciicircum{}m)} with \texttt{m + 1 = k} (ℤ-indexed, no ℕ subtraction) and \texttt{H\textasciicircum{}\{0,k\}\_∂̄(X) :=} that for \texttt{dolbeault E X} (\texttt{d = dbarL}). \texttt{K4 : ((⟨unitCohesive (dolbeault E X) rfl⟩ : HoP \_) ⟶ (⟨…⟩ : HoP \_)⟦k⟧) ≃ₗ[ℂ] H\textasciicircum{}\{0,k\}\_∂̄(X)}. Sanity: both sides are \texttt{0} for \texttt{k < 0} and \texttt{k > n}. For \texttt{k = 0} it is \texttt{O(X)} (F12).
\end{itemize}
\paragraph{5. Items, files, estimates (new files only; no FROZEN changes)}
\begin{itemize}
\item Item: Z2.a; File: \texttt{Algebra/Shift.lean}; Content: \texttt{ι\_a}, generic transport lemma, \texttt{CohesiveModule.shift}; Est. lines: 250–350
\item Item: Z2.b; File: \texttt{Algebra/ShiftHoP.lean}; Content: shift on \texttt{Hom}/\texttt{Z0}/\texttt{H0}, \texttt{HasShift (HoP 𝒜) ℤ} (\texttt{hasShiftMk}), additivity, \texttt{k}-linearity; Est. lines: 300–450
\item Item: K4.a; File: \texttt{Algebra/HomShift.lean}; Content: \texttt{Hom\_\{HoP\}(M, N⟦k⟧) ≃ₗ H\textasciicircum{}k(Hom\textasciicircum{}•(M,N))}; Est. lines: 100–150
\item Item: K4.b; File: \texttt{Algebra/UnitModule.lean}, \texttt{ChecksMain}/\texttt{ChecksForms}; Content: unit module, \texttt{Hom\textasciicircum{}•(U,U) ≅ (𝒜, d)}, Dolbeault cohomology, K4; Est. lines: 200–350
\item Item: Z2.c; File: \texttt{Algebra/TotProd.lean}, \texttt{Algebra/HoPZero.lean}; Content: \texttt{Tot.prodEquiv} (v1.18, below), zero object; Est. lines: 100–200
\item Item: Z2.d; File: \texttt{Algebra/Cone.lean}; Content: cone, structure maps, identities, \texttt{δ} formulas; Est. lines: 350–500
\item Item: Z2.e; File: \texttt{Algebra/ConeMaps.lean}; Content: \texttt{lift}/\texttt{desc}, maps of cones from homotopies; Est. lines: 200–300
\item Item: Z2.f; File: \texttt{Algebra/Pretriangulated.lean}; Content: the pretriangulated axioms, shift of cones; Est. lines: 400–600
\item Item: Z2.g; File: \texttt{Algebra/Triangulated.lean}; Content: octahedron, \texttt{IsTriangulated}; Est. lines: 250–400
\end{itemize}
Total: Z2 core ≈ 1,800–2,700 lines, K4 ≈ 300–500 lines. Order: Z2.a → Z2.b → K4, which needs no cones, then Z2.c–g. This can run in parallel with Z5 (disjoint files: Algebra layer vs Analysis and Forms).

Optional, not proposed now: Z2.h, \texttt{J.CommShift} and \texttt{J.IsTriangulated}. With them \texttt{Perf = J.essImage} is triangulated by Mathlib's \texttt{essImage} instance (+400–700 lines).

\paragraph{6. Risks}
\begin{itemize}
\item **Coherence of \texttt{HasShift}** (\texttt{M[a][b] ≅ M[a+b]}, \texttt{zero\_add}, \texttt{assoc}). This is dependent-type reindexing of \texttt{Q (p + a + b)} vs \texttt{Q (p + (a + b))}. Use a \texttt{reindex (h : f = g)} family with \texttt{subst}, as Mathlib's \texttt{XIsoOfEq}.
\item \textbf{Cone generators.} \texttt{⨁ (Q\_N × Q\_M[1])} is not a type-level biproduct of the \texttt{Tot}s, so the five identities need generator-level proofs (\texttt{homEquivGen}).
\item \textbf{Base-change elaboration traps} (PROGRESS) do not arise here; \texttt{HoP} is over a fixed \texttt{𝒜}.
\end{itemize}
\paragraph{7. Lean targets (to be compiled in \texttt{Scratch/P14\_TR.lean} before skeletons)}
\begin{verbatim}
def CohesiveModule.shift (M : CohesiveModule 𝒜) (a : ℤ) : CohesiveModule 𝒜
instance : HasShift (HoP 𝒜) ℤ
instance : HasZeroObject (HoP 𝒜)
instance (n : ℤ) : (shiftFunctor (HoP 𝒜) n).Additive
instance : Pretriangulated (HoP 𝒜)
instance : IsTriangulated (HoP 𝒜)
def homShiftEquiv (M N : HoP 𝒜) (k : ℤ) : (M ⟶ N⟦k⟧) ≃ₗ[k] HomCohomology M.as N.as k
def dolbeaultCohomology (k : ℤ) : Type
theorem K4 (k : ℤ) : Nonempty ((unitHoP (dolbeault E X) ⟶ (unitHoP (dolbeault E X))⟦k⟧) ≃ₗ[ℂ]
    dolbeaultCohomology E X k)
\end{verbatim}
\paragraph{Blueprint notes.}
octahedron, \texttt{IsTriangulated}

Files: \texttt{Algebra/Triangulated.lean}.
\end{proof}
